% ============================================================================
% SEC05.TEX - Section 5.5: Player Movement
% ============================================================================

\section{Player Movement}

\begin{learningoutcomes}
\item Mengimplementasikan player movement 2D dan 3D
\item Menggunakan Character Controller vs Rigidbody
\item Membuat movement yang smooth dan responsif
\end{learningoutcomes}

\subsection{Character Controller Movement}

\begin{lstlisting}[language={[Sharp]C}]
using UnityEngine;

public class PlayerMovement : MonoBehaviour
{
    public float speed = 5f;
    public float jumpHeight = 2f;
    private CharacterController controller;
    private Vector3 velocity;
    private bool isGrounded;
    
    void Start()
    {
        controller = GetComponent<CharacterController>();
    }
    
    void Update()
    {
        isGrounded = controller.isGrounded;
        
        if (isGrounded && velocity.y < 0)
        {
            velocity.y = -2f;
        }
        
        // Input
        float horizontal = Input.GetAxis("Horizontal");
        float vertical = Input.GetAxis("Vertical");
        
        Vector3 move = transform.right * horizontal + transform.forward * vertical;
        controller.Move(move * speed * Time.deltaTime);
        
        // Jump
        if (Input.GetButtonDown("Jump") && isGrounded)
        {
            velocity.y = Mathf.Sqrt(jumpHeight * -2f * Physics.gravity.y);
        }
        
        // Apply gravity
        velocity.y += Physics.gravity.y * Time.deltaTime;
        controller.Move(velocity * Time.deltaTime);
    }
}
\end{lstlisting}

\begin{figure}[ht]
    \centering
    \begin{tikzpicture}[
        node distance=2cm,
        auto,
        block/.style={rectangle, draw, fill=blue!10, text width=2.5cm, text centered, rounded corners, minimum height=1.2em},
        line/.style={draw, -latex'}
    ]
        % Character Controller Flow
        \node [block] (input1) {Input};
        \node [block, right=of input1] (code1) {Code (Move)};
        \node [block, right=of code1] (transform1) {Transform Position};
        \node [block, above=0.2cm of code1, draw=none, fill=none] {Character Controller};
        
        \path [line] (input1) -- (code1);
        \path [line] (code1) -- (transform1);
        
        % Rigidbody Flow
        \node [block, below=1.5cm of input1] (input2) {Input};
        \node [block, right=of input2] (force) {AddForce / Velocity};
        \node [block, right=of force] (physics) {Physics Engine};
        \node [block, right=of physics] (transform2) {Transform};
        \node [block, above=0.2cm of force, draw=none, fill=none] {Rigidbody};
        
        \path [line] (input2) -- (force);
        \path [line] (force) -- (physics);
        \path [line] (physics) -- (transform2);
    \end{tikzpicture}
    \caption{Perbandingan Alur Character Controller vs Rigidbody}
    \label{fig:movement_comparison}
\end{figure}

\begin{catatan}
Character Controller memberikan kontrol penuh atas movement tanpa fisika. Rigidbody menggunakan fisika Unity untuk movement yang lebih realistis.
\end{catatan}
